%!TEX TS-program = xelatex
% vim: set fenc=utf-8

% -*- coding: UTF-8; -*-
%!TEX encoding = UTF-8
\documentclass[twoside]{CUGThesis}

\title{论文题目} %论文题目
\author{作者姓名} %作者姓名
\date{\today} %日期，默认当日
\school{院系名称} %院系名称
\classnum{专业} %专业
\stunum {学号} %学号
\instructorone{某某某} %指导教师1姓名
\instructoronelevel{副教授}
\instructortwo{某某} %指导教师2姓名
\instructortwolevel{副教授}

\usepackage{makecell}
\usepackage{multicol}
\usepackage{multirow}
\usepackage{hyperref}

\begin{document}
	\maketitle
	\makestatement
	
	\begin{cnabstract}{关键词；关键词；关键词；关键词}
		摘要
	\end{cnabstract}
	
	\begin{enabstract}{keyword; keyword; keyword; keyword}
		abstract
	\end{enabstract}
	
	\makeToc
	
%---------------------------------------------开始正文---------------------------------------------
	
%--------------------------------------------- 第一章 ---------------------------------------------
	\begin{spacing}{2}
		\section{第一章}
	\end{spacing}
	引言
	\subsection{第一章第一小节}
	第一小节，引用实例\cite{引用}，按照格式填写mybib.bib
	\subsection{第一章第二小节}
	第二小节
	\subsection{etc}
	etc
	
	
%------------------------------------------- 第二章 ---------------------------------------------
	\begin{spacing}{2}
		\section{第二章}
	\end{spacing}
	第二章引言
	\subsection{第二章第一小节}
	第二章第一小节
	\subsubsection{第二章第一小节第三级标题}
	插图实例（图\ref{Fig:example}）
	\begin{figure}[!t]
		\centering
		\includegraphics[scale=0.5]{Figures/example.jpg}
		\caption{示例图}
		\label{Fig:example}
	\end{figure}
	\subsection{第二章第二小节}
\begin{table}[t]
	\centering
	\caption{在最短路问题中PL和ACO算法的比较}
	\label{Tab:ACOcomparison}
	\scalebox{0.50}{
		\begin{tabular}{ccc}
		\toprule  %添加表格头部粗线
		姓名& 学号& 性别\\
		\midrule  %添加表格中横线
		Steve Jobs& 001& Male\\
		Bill Gates& 002& Female\\
		\hline \hline
		\end{tabular}
	}
	\vspace{10mm}
\end{table}

	
%--------------------------------------------- 第三章 ---------------------------------------------
	\begin{spacing}{2}
		\section{第三章}
	\end{spacing}
	引言
	\subsection{第三章第一小节}
	开始列表
	\begin{enumerate}
		\item 条目1
		
		说明文字1
		
		\item 条目2
		
		说明文字2
		
		
	\end{enumerate}
	
	
%---------------------------------------------  致谢  ---------------------------------------------
	\begin{spacing}{2}
		\section*{致谢}
	\end{spacing}
	\phantomsection
	\addcontentsline{toc}{section}{致谢}
	
	致谢
	
	\clearpage
%---------------------------------------------参考文献---------------------------------------------
	
	\bibliography{Bibs/mybib}
	\phantomsection
	\addcontentsline{toc}{section}{参考文献}
	
\end{document}
